% !TeX root = ../深度学习与大语言模型技术原理.tex
% 由主文档通过 \input 载入；请编译主文档。
\setcounter{chapter}{26} % 保留原书章号。
\chapter{视觉、音频、视频与统一多模态模型}
\label{ch:multimodal}

\section{视觉编码与连接器}

多模态大语言模型（multimodal large language model, MLLM）把文本之外的信号转成模型可使用的表示。视觉 Transformer（Vision Transformer, ViT）将图像切成图像块（patch），每块展平后线性投影，加入位置信息再经注意力编码\cite{vit}。若图像大小为 $H\times W$、块边长为 $p$ 且可整除，块数为 $HW/p^2$；提高分辨率会增加视觉序列长度。

视觉编码器输出 $H_v\in\mathbb R^{n_v\times d_v}$，连接器（connector）或投影器（projector）把它映到语言模型维度，例如 $Z=H_vW_p$、$W_p\in\mathbb R^{d_v\times d}$。更复杂的连接器可以使用 MLP、可学习查询或交叉注意力，把大量视觉词元压缩为固定数量。压缩节约上下文，却可能丢失文字、小物体或空间关系。

视觉词元可与文本交错输入，或由语言层通过交叉注意力访问。占位符、图像顺序、分辨率变换和位置编码必须一致；“看见图片”不是一个与预处理无关的能力。多个裁剪块还需说明其原图坐标关系，否则模型难以恢复全局空间结构。

\section{图文对齐与指令微调}

图文对比学习使匹配图像与文本向量相近，非匹配向量相远。对批内 $B$ 对样本，令单位向量为 $v_i,t_i$，相似度 $s_{ij}=v_i^\top t_j/\tau$，可对图到文与文到图分别使用交叉熵并求平均。CLIP 是这种对齐路线的代表\cite{clip}；相似度对齐本身不训练自回归长回答。

视觉指令微调在图像、指令和回答上计算语言目标，可先训练连接器再训练部分或全部模型，也可采用其他联合策略。LLaVA 的 Visual Instruction Tuning 展示了视觉连接器与语言指令数据结合的路径\cite{llava}。冻结视觉编码器能减少训练状态，却限制其适配新视觉领域的能力。

视觉定位将词语关联到图像区域，光学字符识别（optical character recognition, OCR）关注文字内容。对象识别正确不代表相对位置、计数和文字读取都正确。评测要防止语言先验掩盖视觉错误，例如遮挡图像后仍给同一答案，说明模型可能主要依赖文本偏置。

\section{音频与视频的时序结构}

音频可以经波形或时频表示编码，也可以转成离散音频码。语音识别、声音事件理解、说话人信息和语音生成的目标不同；把音频先转写成文本会丢失语调和非语言声学信息。实时交互还需管理分块输入、增量输出和被打断后的状态。

视频同时包含空间与时间。稀疏抽帧降低成本，却可能漏掉短暂关键事件；密集抽帧增加大量相似视觉词元。帧时间戳、音视频同步和跨帧身份跟踪影响时序推理，不能把视频简单当成无序图片集合。长视频摘要还需平衡局部细节与全局事件组织。

\section{统一多模态与模态间评测}

统一多模态模型常被称为 Omni Model，但这一名称不对应唯一架构。系统可以使用不同模态编码器、统一语言主干和多个输出头，也可采用更统一的词元表示。需要说明支持的是输入理解、输出生成还是实时双向交互，不能仅凭“多模态”断言全链路能力。

交错数据（interleaved data）包含文本、图像、音频等交替片段，训练时必须对齐模态边界、时间与目标掩码。不同模态损失尺度和数据量差异很大，混合比例会影响训练方向。评测应覆盖单模态能力、跨模态绑定、缺失模态、噪声与冲突信息，并单独报告感知错误和推理错误。

